% ============================================================
%  Chapter 4: R4 — Quantum Postulates
% ============================================================
\chapter{R4 — 量子力学公设}
\label{ch:r4}

\section{总说：世界的真正语言}

R4 是整个物理学中最具革命性的原理。它说：\textbf{物理世界的基本描述语言不是确定性的轨道，而是 Hilbert 空间中的态向量}。

量子力学包含四条相互关联的公设。它们共同定义了从经典物理到量子物理的过渡——这不是一个渐进变化，而是一个\textbf{结构性分界}。

\begin{principlebox}[R4 — 量子力学公设]
  \textbf{(4a) 叠加原理}：合法量子态的线性组合仍为合法量子态。

  \textbf{(4b) 幺正演化}：孤立系统的演化由幺正算符实现，保证概率守恒。

  \textbf{(4c) 正则对易关系}：$[\hat{x}_i, \hat{p}_j] = i\hbar\delta_{ij}$——量子与经典的结构分界线。

  \textbf{(4d) 算符-可观测量对应}：每个物理可观测对应一个厄米算符；测量结果必为其本征值。

  \textbf{注}：R4d 仅声明"测量结果 = 本征值"，不包含投影公设。R4+R5 构成\textbf{最小统计诠释}——对测量问题给出完备概率预言，不承诺坍缩本体论。
\end{principlebox}

% ── 变量定义表 ──
\section{变量定义}

\begin{table}[h]
\centering
\small
\begin{tabular}{p{2.4cm}p{1.8cm}p{2.4cm}p{1.8cm}p{3cm}p{3cm}}
\toprule
\textbf{变量} & \textbf{符号} & \textbf{类型} & \textbf{量纲} & \textbf{定义域} & \textbf{物理含义} \\
\midrule
量子态 & $|\psi\rangle$ & $\mathcal{H}$ 中向量 & [1] & $\||\psi\rangle\| = 1$ & 系统全部信息的完备编码 \\
Hilbert 空间 & $\mathcal{H}$ & 复内积空间 & — & $\dim\mathcal{H} \ge 1$ & 量子态的可能取值空间 \\
基态 & $|i\rangle$ & $\mathcal{H}$ 中向量 & [1] & $\langle i|j\rangle = \delta_{ij}$ & 完备正交归一基 \\
概率幅 & $\alpha_i$ & $\mathbb{C}$ & [1] & $\sum_i |\alpha_i|^2 = 1$ & 态在基态 $|i\rangle$ 的投影 \\
幺正演化算符 & $U(t_2,t_1)$ & $\mathcal{H}\to\mathcal{H}$ & [1] & $U^\dagger U = UU^\dagger = I$ & 时间平移；保概率守恒 \\
哈密顿算符 & $\hat{H}$ & 厄米算符 & M$\cdot$L$^2\cdot$T$^{-2}$ & $H^\dagger = H$ & 能量算符；演化生成元 \\
约化 Planck 常数 & $\hbar$ & 常数 & M$\cdot$L$^2\cdot$T$^{-1}$ & $\hbar = 1.055\times10^{-34}$ J$\cdot$s & 量子作用量单元 \\
位置算符 & $\hat{x}$ & 厄米算符 & L & $\hat{x}^\dagger = \hat{x}$ & 位置可观测量 \\
动量算符 & $\hat{p}$ & 厄米算符 & M$\cdot$L$\cdot$T$^{-1}$ & $\hat{p}^\dagger = \hat{p}$ & 动量可观测量 \\
对易子 & $[A,B]$ & 算符 & 依赖 $A,B$ & $[A,B] \equiv AB-BA$ & 是否可同时对角化 \\
可观测量算符 & $\hat{O}$ & 厄米算符 & 依赖量 & $\hat{O}^\dagger = \hat{O}$ & 物理可观测量的表示 \\
本征值 & $o_i$ & $\mathbb{R}$ & 依赖 $\hat{O}$ & $\hat{O}|o_i\rangle = o_i|o_i\rangle$ & 可能测量结果 \\
本征态 & $|o_i\rangle$ & $\mathcal{H}$ 中向量 & [1] & $\hat{O}|o_i\rangle = o_i|o_i\rangle$ & 确定值态 \\
\bottomrule
\end{tabular}
\caption{R4 变量定义}
\end{table}

% ── 数学表达式 ──
\section{数学表达式}

\subsection{(4a) 叠加原理}
\[
|\psi\rangle = \sum_{i} \alpha_i |i\rangle, \quad \sum_i |\alpha_i|^2 = 1
\]
$\{|i\rangle\}$ 为 $\mathcal{H}$ 的完备正交归一基：$\langle i|j\rangle = \delta_{ij}$。

\subsection{(4b) 幺正演化}
\[
|\psi(t)\rangle = U(t, t_0) |\psi(t_0)\rangle, \quad U^\dagger U = I
\]
\[
U(t, t_0) = \exp\!\left(-\frac{i}{\hbar} \hat{H} (t - t_0)\right)
\]
微分形式：$i\hbar \frac{d}{dt}|\psi(t)\rangle = \hat{H} |\psi(t)\rangle$（抽象的 Hilbert 空间演化方程，\textit{非}坐标表象 Schrödinger 方程）。

\subsection{(4c) 正则对易关系}
\[
[\hat{x}_i, \hat{p}_j] = i\hbar\,\delta_{ij}
\]
坐标表象：$\hat{p} \equiv -i\hbar\nabla$。

\subsection{(4d) 算符-可观测量对应}
\[
\hat{O}|o_i\rangle = o_i|o_i\rangle, \quad o_i \in \mathbb{R}
\]
本征态 $\{|o_i\rangle\}$ 构成完备正交归一基。

% ── 成立条件与失效边界 ──
\section{成立条件与失效边界}

\begin{limitbox}[R4 的成立条件与失效边界]
  \textbf{成立条件}：
  \begin{itemize}
    \item 量子区 (quantum\_regime)：作用量 $\sim \hbar$ 时量子效应显著；$\hbar \to 0$ 恢复经典
    \item 孤立系统 (isolated\_system)：R4b 仅适用于未测量的孤立系统
    \item 注：R4a-d 本身框架无关——不预设相对论。非相对论极限仅在 $\hat{H} = -\hbar^2\nabla^2/2m + V$ 时需要
  \end{itemize}

  \textbf{失效边界}：
  \begin{itemize}
    \item 测量过程：幺正演化中断；R5 接管
    \item Planck 尺度量子引力：时间可能不再是参数
    \item 开放量子系统：非幺正演化需 Lindblad 方程
  \end{itemize}
\end{limitbox}

% ── 必要性 ──
\section{必要性 (Necessity)}

\begin{table}[h]
\centering
\small
\begin{tabular}{p{5cm}p{7cm}}
\toprule
\textbf{移除/弱化} & \textbf{后果} \\
\midrule
无叠加原理 & 无干涉、无纠缠 $\to$ 退化为经典概率论 \\
无幺正演化 & 概率不守恒；无稳定量子态 \\
$[\hat{x},\hat{p}] = 0$（经典极限） & 位置动量可同时确定 $\to$ 无不确定性 $\to$ 无离散能谱、无隧穿 \\
无可观测量-算符对应 & 无法从态向量提取物理预言 \\
\bottomrule
\end{tabular}
\caption{R4 必要性}
\end{table}

% ── 充分性 ──
\section{充分性 (Sufficiency)}

\textbf{R4 + R1} 提供：
\begin{itemize}
  \item 单粒子量子力学（Schrödinger 方程、能级量子化、隧穿）
  \item + R2 $\to$ 路径积分量子化
  \item + R1 相对论推广 $\to$ 量子场论（Klein-Gordon, Dirac, Yang-Mills）
  \item 自旋-统计定理（R4 + R1 $\to$ Pauli 不相容原理）
\end{itemize}
